\par\noindent\begin{minipage}{\linewidth}
% !TeX root = ../教材题库_全部条目_转写工作稿.tex
% Source: XJ-B1-C01-SU-P028-Y1404
% 第1章 / 小结与复习 小结与复习 / 教材 p.28 / PDF 第35页
% 已对照原页转写题干及公式；待教师复核；未制作答案
% 原题及相邻原解截图（核对用）：../题目原图/XJ-B1-C01-SU-P028-Y1404-01.webp
% 本文件只在转写完成后作为题干使用；OCR 条目尚不是可靠题干。
\begin{bankitem}{第1章 小结与复习 第5题}{第1章 / 小结与复习 小结与复习 / 教材 p.28 / PDF 第35页}{XJ-B1-C01-SU-P028-Y1404}{已对照原页转写题干及公式；待教师复核；未制作答案}
数学中的大量命题具有“若 $p$，则 $q$”的形式，试通过对这些命题的梳理，说明充分条件、必要条件以及充要条件的意义。
\end{bankitem}
\end{minipage}\par\medskip
